\scanpage{163}
Si $\widetilde\xi$ est réelle, on a une structure de Hodge à involution.

Comment se relient ces structures de Hodge ??

\diagram{\begin{tikzcd}[column sep=huge]
\substack{\widetilde{\mathcal M}_0\\\text{motifs}\\\text{abéliens sur}\\\widetilde K\ \text{provenant}\\\text{de }K} \arrow[r,shift left=5,"{\text{fibre en }\widetilde\xi}"] \arrow[r,shift right=5,"{\text{fibre en }\widetilde\xi'}"'] & \deleted{\text{\unsure{Vecteurs}}}\ G^0\text{-Modules} \arrow[r,shift left=5,"\Phi_{\widetilde\xi}"] \arrow[r,shift right=5,no head,"\unsure{\Phi|\varphi}"'] & \substack{\text{Réseaux de}\\\text{Hodge}}
\end{tikzcd}}

\deleted{\unsure{NB} Le \unsure{degré} de $\xi$}

\deleted{\uline{NB} Le seul \deleted{\unsure{bon}} caractère \unsure{invariant|rationnel} de $\uuline{K}^*$ est $N_{K^*/\mathbf Q}$ ; ainsi \deleted{\unsure{les}}
\diagram{\begin{tikzcd}[column sep=large]
\uuline{K}^*\text{-modules} \arrow[r,"\sim"] & \substack{\text{\unsure{systèmes} de}\\\text{vectoriels de dim finie sur }K}
\end{tikzcd}}
\deleted{\unsure{modules}} Si $E^0\ne0$, l'opération est fidèle pour $E^0$ ; le \unsure{sous-espace} est nul. Donc les vectoriels sur lesquels $\uuline{K}^*$ opère sont les \deleted{vectoriels} espaces \inserted{vectoriels} gradués par $\mathbf Z$.}

\deleted{$V$ un $G^0$-module

$V\otimes_{\mathbf Q}\mathbf C$ muni de \deleted{une} structure de Hodge --- $\unsure{\pi_\xi^V|\mathcal H_\xi^V}$

($\xi$ place de $\widetilde K$)

et d'un sous-$\widetilde K$-espace vectoriel $H_\xi^V$}
\margin{162}
